\documentclass[10pt,a4paper]{amsart}
\usepackage[margin=19mm]{geometry}
\usepackage{amsmath,amssymb,mathtools}
\usepackage[scr=rsfs,cal=euler]{mathalfa}
\usepackage{xfrac}
\usepackage[unicode,hidelinks,bookmarks=false]{hyperref}
\newcommand{\ie}{i.e.\ }
\newcommand{\cf}{cf.\ }
\newcommand{\resp}{respectively\ }
\newcommand{\Subsection}{\subsection}
\providecommand{\nobreakdash}{\nobreak\hbox{-}}
\setlength{\emergencystretch}{2em}
\multlinegap0pt
\usepackage{lipsum}
\usepackage{epstopdf}
\usepackage{algorithmic}
\usepackage{color}
\usepackage{enumitem}
\newcommand{\creflastconjunction}{, and~}
\newtheorem{theorem}{Theorem}[section]
\newtheorem{lemma}[theorem]{Lemma}
\newtheorem{corollary}[theorem]{Corollary}
\newtheorem*{claim}{Claim}
\newtheorem{definition}[theorem]{Definition}
\newtheorem{remark}[theorem]{Remark}
\newcommand{\R}{\mathbb{R}}
\newcommand{\ep}{\epsilon}
\newcommand{\F}{\mathcal{F}}
\newcommand{\N}{\mathcal{N}}
\newcommand{\C}{\widetilde{C}}
\newcommand{\E}{\mathcal{E}}
\newcommand{\G}{\mathcal{G}}
\newcommand{\M}{\mathcal{M}}
\newcommand{\PA}{\mathcal{P}}
\newcommand{\SA}{\mathcal{S}}
\newcommand{\X}{\mathcal{X}}
\newcommand{\0}{\textbf{0}}
\newcommand{\HA}{\mathcal{H}}
\usepackage{amsopn}
\DeclareMathOperator{\diag}{diag}
\begin{document}
\title[Exploding and vanishing gradients in deep neural networks: t]{Exploding and vanishing gradients in deep neural networks: the effect of residual connections}
\author{Vivek S Borkar}
\date{}
\maketitle
\noindent\emph{Source and licence.} Exploding and vanishing gradients in deep neural networks: the effect of residual connections. Version arXiv:2606.17013v1 (CC BY 4.0). This material is distributed under the Creative Commons Attribution 4.0 International licence, http://creativecommons.org/licenses/by/4.0/, as recorded in the arXiv OAI licence metadata for this exact version and shown on the abstract page https://arxiv.org/abs/2606.17013v1. Full rights notice: licenses/arxiv-2606.17013-CC-BY-4.0.txt.\par\smallskip
\noindent\textit{Editorial scope.} Counted unit: the original \texttt{theorem} environment \texttt{final} at source lines 325--327 together with its complete proof at lines 329--338; the delivered document keeps source lines 151--339 so that every referenced label and every citation resolves inside it. Native editable source is the paper's own concatenated TeX; the AMS wrapper is editorial formatting only. No publisher class, upstream script or unrelated artwork is redistributed.
\begin{equation}
f^N_{\beta^N} := f_{\beta_N}\circ f_{\beta_{N-1}} \circ \cdots \circ f_{\beta_1} \circ f_{\beta_0} : \R^d \to \R^d. \label{DNNmap}
\end{equation}
Here `$\circ$' denotes composition of functions, i.e.\ $f\circ g(x) := f(g(x))$. \\

There is a standard way of mapping this composition into iterates of an equivalent discrete time dynamical system. Specifically, let $x_0 \in \R^d$ and recursively define 
\begin{equation}
x_{n+1} = f_{\beta_n}(x_n), \ n \geq 0. \label{dynamics0}
\end{equation}
Then inductively, one has
$$f^N_{\beta^N}(x_0) = x_N, N \geq 1$$
and vice versa. Thus the two descriptions are completely equivalent.\\

Let  $D(\beta)(x) :=$  the $d\times d$ Jacobian matrix of $f_\beta$ evaluated at $x$, parametrized by $\beta$. Likewise, let $D^N(\beta^N)(x^N):=$  the $d\times d$ Jacobian matrix of  $f^N$, parametrized by $\beta^N$ and evaluated at $x^N := [x_N, \cdots , x_0]$. Then by the chain rule of differentiation,
\begin{equation}
D^N(\beta^N)(x^N) = \prod_{m=0}^{N}D(\beta_{N-m})(x_{N-m}). \label{chainrule}
\end{equation}
Let $\eta_0\in \R^d$. We shall be interested in the asymptotic behaviour of 
\begin{eqnarray}
\eta^N &:=& \left(\prod_{m=0}^{N}D(\beta_{N-m})(x_{N-m})\right)\eta_0. \label{mult} \\
&=& D(\beta_N)(x_N)\eta^{N-1}. \label{dynamics1}
\end{eqnarray}
Equation \eqref{dynamics1} represents a time-inhomogeneous linear dynamical system with $\beta_n$'s serving as parameters. The second equality is the counterpart of the equivalence of  \eqref{DNNmap} and \eqref{dynamics0} mentioned above, but for the \textit{linearized dynamics} given by \eqref{dynamics1} as opposed to the original dynamics \eqref{dynamics0}.\\

Summarizing, this dynamical view of DNNs is tantamount to:\\

\begin{enumerate}
\item viewing input as the initial condition of a nonlinear dynamics,\\

\item the layer index $m \geq 0$ as a discrete time index,\\

\item the input-output map $f_{\beta_i}$ of the $i$th layer of the feedforward network as a time-dependent map that dictates the evolution of the dynamics at time $i$ from the state at time $i$ to the state at time $i+1$,\\

\item $\beta_i$'s are parameters of the next step map at time $i$, $i \geq 0$. 
 
 \end{enumerate}
 
 \medskip
 
 This equivalence allows us to apply techniques from the theory of nonlinear dynamics to DNNs as the depth of the DNN tends to infinity.\\
 
 \subsection{Notation}

 
 We shall use the following additional notation.\\
 
 \begin{enumerate}

\item We denote by $\M$ the set of $d\times d$ matrices and by $\M^+ \subset \M$ the subset thereof consisting of nonsingular $d\times d$ matrices. \\

\item We denote by $\0$ the zero vector in the appropriate dimension, depending on the context.\\

\item We denote by $\PA(\X)$ the Polish\footnote{i.e., a separable Hausdorff space which has a compatible complete metric} space of probability measures on the Polish space $\X$ with Prokhorov topology (also known as the topology of weak convergence). This topology is defined as the coarsest topology that renders continuous the maps $\zeta \in \PA(\X) \mapsto \int fd\zeta\in \R$ for $f \in C_b(\X) :=$ the space of bounded continuous functions on $\X$. (See, e.g., \cite{Bill} for a detailed exposition of this topology.)\\

\item We denote by $S^d$ the unit sphere in $\R^d$ and denote by $\Phi : \R^d\backslash\{\0\} \to S^d$ the map $x \mapsto \frac{x}{\|x\|}$ when $x \neq \0$.\\

\item We denote by $P^{d-1}$ the projective space, i.e., the space of  equivalence classes of nonzero vectors in $\R^d$ under the equivalence relation $x \equiv \lambda x$ for $\lambda \in \R\backslash\{\0\}$. If $u \in P^{d-1}$, we denote by $\hat{u}$ a generic element of $\R^d$ that gets mapped to $u$ under this equivalence relation. \\

\item We denote by $\Psi : \R^d\backslash\{\0\}  \to P^{d-1}$ the map that maps $x \in \R^d$ to the corresponding element of $P^{d-1}$. $P^{d-1}$ is endowed with the natural metric 

$$d(u,u') \ := \ \min\left(\left\|\frac{\hat{u}}{\|\hat{u}\|}-\frac{\hat{u}'}{\|\hat{u}'\|}\right\| \ , \ \left\|\frac{\hat{u}}{\|\hat{u}\|}+\frac{\hat{u}'}{\|\hat{u}'}\right\|\right),$$

\medskip

for $\hat{u}\in \Psi^{-1}(u)$ and $\hat{u}' \in \Psi^{-1}(u')$.\\

\item $\mu$ is a prescribed probability measure on $\M$.

\end{enumerate}

\medskip

We next recall the Furstenberg-Kifer theorem of multiplicative ergodic theory, which will be the basis of our analysis of DNNs.


\section{Furstenberg-Kifer theorem}\label{main}

Multiplicative ergodic theorems go back to \cite{Osel}. The basic result and the associated theory of Liapunov exponents has been extended in many directions, see \cite{Arnold} for a comprehensive treatment. We use here a variant due to Furstenberg and Kifer \cite{FK} which has the additional feature of giving a clean characterization of the Liapunov spectrum and the resulting direct sum decomposition of the state space, in terms of an associated Markov chain on the projective space. This needs the following additional assumption.\\

\noindent \textbf{Assumption 1:}\label{Ass} $$\int\left(\log^+\|A\| + \log^+\|A^{-1}\|\right)\mu(dA) < \infty.$$

\medskip


%This immediately implies the following.
%
%\begin{lemma}  $D^N( \cdot ), D( \cdot )$ above take values in $\M^+$, $\mu$-a.s.. \\
%\end{lemma}

\begin{definition}\label{MKernel} Given a $\mu \in \PA(\M)$ and a $\nu \in \PA(P^{d-1})$, we define $\mu*\nu \in \PA(P^{d-1})$ by
\begin{equation}
\int \varphi d(\mu*\nu) := \int \varphi(Ax)\mu(dA)\nu(dx) \ \ \ \ \forall \ \varphi \in C_b(\M). \label{star}
\end{equation}
\end{definition}

\medskip

\begin{lemma}\label{extreme} The set $\SA := \{\nu \in \PA(P^{d-1}) : \mu*\nu = \nu\} \subset \PA(P^{d-1})$ is a nonempty  and compact simplex whose extreme points are mutually singular. \end{lemma}

\begin{proof} This is immediate from the discussion in \cite{FK}, pp.\ 17-18, where the correspondence of $\SA$ with the stationary distributions of a Markov chain is established. Specifically, define the transition kernel 
$$\kappa : P^{d-1} \to \PA(P^{d-1})$$ 
by
$$\int f(y)\kappa(dy|x) := \int_{P^{d-1}}f\left(\frac{Ax}{\|Ax\|}\right)\mu(dA) \ \ \forall \ f \in C\left(P^{d-1}\right).$$
(Equivalently, $$\kappa(B|x) := \int_{P^{d-1}}I\left\{\left(\frac{Ax}{\|Ax\|}\right)\in B\right\}\mu(dA)$$
for all Borel $B \subset P^{d-1}$.) \\

Thus $\nu$ is simply an invariant measure of this transition kernel by Definition \ref{MKernel}. Since $P^{d-1}$ is compact and the map $x \to \kappa(dy|x)$ is  seen to be continuous in $x$, it follows from standard Markov process theory \cite{Benaim}, \cite{MeynTweedie} that the set of such probability measures $\nu$ forms a nonempty compact simplex whose extreme points are mutually singular. \end{proof}

\medskip

We call $\SA$ the set of $\mu$-stationary measures in $\PA(P^{d-1})$. We shall also need the following definitions.\\

\begin{enumerate}

\item For a subspace $Z$ of $\R^d$, denote by $\overline{Z}$ the set of corresponding elements of $P^{d-1}$.\\

\item For $\mu \in \PA(\M)$, say that a subspace $Z$ of $\R^d$ is  \textit{$\mu$-invariant} if it is invariant under $\mu$-a.s.\ $A \in \M$.
\end{enumerate}

\begin{remark}\label{DIM} A priori, our assumption that all $f_{\beta_i}$'s map $\R^d$ to itself is not restrictive. We can take the $d < \infty$ to be an upper bound (assumed to exist) on the dimensionality of the input or output space of the individual networks and set the appropriate connection weights to zero if the actual dimension is lower. However,  Assumption 1 above is  restrictive. It is required in order to be able to use the theory of \cite{FK}.  Our aim is to demonstrate a plausible mechanism for explaining the observed benefits of residual connections in a quantitative fashion, albeit for a stylized model. Since matrices satisfying Assumption 1 are dense in $\M$, this exercise is not entirely unreasonable. That said, it will need a lot more sophisticated mathematics (in particular, a suitable extension of Theorem \ref{FuKe} in order to push these results to full generality, which is a task for the future. \end{remark}

\medskip

With this notation and caveats, we next state the key results from \cite{FK} that are relevant for our purposes. These have been recast in our notation.\\

Under stated hypotheses, \cite{FK} proves the following (See Theorems 3.9 and 3.10 of \textit{ibid.}).

\begin{theorem}\label{FuKe} $(i)$ There exist an integer $1 \leq r \leq d$, a sequences of subspaces of $\R^d$
$$\{0\} \subset L_r \subset L_{r-1} \subset \cdots \subset L_2 \subset L_1 \subset L_0 = \R^d$$
and a sequence of real numbers
$$\gamma(\mu) := \gamma^0(\mu) > \gamma^1(\mu) > \gamma^2(\mu) > \cdots > \gamma^r(\mu)$$
such that, if $v\in L_i\backslash L_{i+1}$, then for $\{X_n\}$ i.i.d.\ with law $\mu$,
$$\lim_{N\uparrow\infty}\frac{1}{N}\log\left\|X_NX_{N-1}\cdots X_1v\right\| = \gamma^i(\mu).$$

$(ii)$ These $\{\gamma^i(\mu), 0 \leq i \leq r\}$ are precisely the discrete values taken by the quantity
\begin{equation}
\alpha(\mu,\nu) := \int_\M\int_{P^{d-1}}\log\left(\|Au\|\right)\mu(dA)\nu(du) \label{alpha}
\end{equation}
as $\nu$ varies over all $\mu$-stationary measures. Also, $L_i :=$ the unique maximal $\mu$-invariant subspace of the set of all $\mu$-invariant subspaces satisfying $\nu(\overline{L}_i) = 0$ for all $\nu$ with $\alpha(\mu,\nu) > \gamma^i(\mu)$.
\end{theorem}

The $\{\gamma^i_\mu\}$ are called the Liapunov spectrum associated with the i.i.d.\ matrices $\{X_n\}$. The significant part of Theorem \ref{FuKe} for our purpose is part $(ii)$, which characterizes the Liapunov spectrum and the associated direct sum decomposition of the state space in terms of a Markov chain.\\

\section{Effect of residual connections}\label{resnet}

In order to map our problem to this framework, we consider $\{\beta_n\}$ i.i.d.\ with law (say) $\Gamma \in \PA(\R^s)$. We make the following additional assumption:\\

\noindent \textbf{Assumption 2:} The limit $x_\infty := \lim_{n\uparrow\infty}x_n$ exists a.s.\ in \eqref{dynamics0}.\\

We justify this as follows. In classifier DNNs, there are finitely many classes encoded as euclidean vectors $C = \{c_1, \cdots , c_k\} \subset \R^d$ (say) at the output and for $\mu$-a.s. $x_0$, the above limit $x_\infty \in C$ is well defined. Fix $x_0 = v$ in the probability $1$ set where this holds. Then Assumption 2 holds.\\

As we are considering the asymptotic regime, in what follows, we consider $x_n$ replaced by $x_\infty$. Furthermore, since $x_\infty$ is measurable with respect to the tail $\sigma$-field $\cap_{n \geq 1}\sigma(\beta_i, i \geq n)$ which is trivial by the Kolmogorov $0-1$ law, $x_\infty$ is a.s. a constant and we may take it to be a deterministic constant. Thus $D(\beta_n)(x_\infty), n \geq 0,$ are i.i.d.\ with law (say) $\upsilon$.\\

The ResNet architecture can be viewed as adding a forward connection, the so called `residual connection', that replaces every forward block $x \mapsto f_{\beta_i}(x)$ of the DNN by the map $x \mapsto x + f_{\beta_i}(x)$. We now explore how this affects the overall input-output map of the DNN in the limit as $N \uparrow \infty$, in view of the foregoing.\\

We do this by comparing the above maps for $x = \hat{u}$ as defined earlier with $\|\hat{u}\| = c > 0$ (say).  Since $\mu$ and therefore the set $\SA$ of possible values of $\nu$ is fixed, we focus on the quantity 
$$\xi(u, A) := \log(\|Au\|) = \log\left(\frac{\|A\hat{u}\|}{\|\hat{u}\|}\right),$$
and explore how it changes when $A$ is replaced by $I + A$.  \\

%Let $S_1, S_2$ be compact Polish spaces and $\HA: S_1 \to S_2$ a homeomorphism between them. Let $\HA_*$ denote the corresponding induced homeomorphism between $\PA(S_1)$ and $\PA(S_2)$, i.e., $\HA_*(\nu) = \hat{\nu}$ where $\hat{\nu}$ is defined by
%$$\int f d\hat{\nu} := \int f(f\circ d\nu \ \ \forall \ f \in C_b(S_1).$$ 
%Likewise we define $(\HA\times\HA)_* : \PA(S_1\times S_1) \to \PA(S_2\times S_2)$. Let $\upsilon \in \PA(S_1\times S_1)$ disintegrate as $\upsilon(ds\times ds') = \upsilon_0(ds)\upsilon_1(ds'|ds)$. \\
%
%\begin{lemma}\label{pushf} For $\HA_*$ defined as above, 
%$$\HA_*(\upsilon_1(dy|x))\HA_*(\upsilon_0)(dx) = (\HA\times\HA)_*(\upsilon(dx,dy)).$$
%Furthermore, if $S_1 = S_2$ and $\upsilon_0(dx)$ is invariant under the transition kernel $\upsilon_1(dy|x)$, then $\HA_*(\upsilon_0)$ is invariant under the transition kernel $\HA_*(\upsilon_1(dy|x))$.
%\end{lemma}
%
%\medskip
%
%This follows by direct verification. 

Now we are ready to prove our main result. \\


\begin{theorem}\label{final}  
The Liapunov spectrum under a residual connection is a smaller perturbation of the spectrum of the identity matrix (i.e., the vector of all $1$'s) than without a residual connection.
\end{theorem}

\begin{proof} It is easy to check that $\Phi(\hat{u}), \Psi(\hat{u})$ are independent of the $c = \|\hat{u}\|$ above from their very definition. Hence we can and do take $c = 1$ without any loss of generality. Now consider  
$$\hat{u}_0 := \hat{u}, \ \hat{u}_1 := A\hat{u}_0, \ \hat{u}_2 := (I + A)\hat{u}_0 = \hat{u}_0 + A\hat{u}_0.$$ 
Consider the two dimensional parallelopiped $B$ formed by $\0, \hat{u}_0, \hat{u}_1$ and $\hat{u}_2$. Let $\hat{u}_3:=$ the intersection of its diagonals. Then $\hat{u}_3$ is in the relative interior of $B$. Let $\hat{u}_i' := \Phi(\hat{u}_i), i = 0,1,3$. Then $u_i' \in B\cap S^d$ for $ i = 0,1,3$. It is then easy to see that $\hat{u}_3'$ lies in the relative interior of the arc $B\cap S^d$ joining $\hat{u}_0'$ and $\hat{u}_1'$. Therefore it is closer to $\hat{u}_0'$ than $\hat{u}_1'$. This property is  preserved under the map $\Psi$. That is, on mapping these vectors to the corresponding equivalence classes in $P^{d-1}$, $\Psi(\hat{u}_3')$ lies closer to $\Psi(\hat{u}_0)$ than $\Psi(\hat{u}_1)$, in the metric topology of $P^{d-1}$. 
%Now in Lemma \ref{pushf}, let $S_1 = S_2 = P^{d-1}$ and $\HA(u) := \Psi((I + A)u)$. 
Hence it follows that
\begin{eqnarray*}
\lefteqn{\int_{P^{d-1}}\int_\M \log\left(\|(I + A)u\|\right))\varphi(dA)\nu(du) \ -} \\
&& \  \int_{P^{d-1}}\int_\M \log\left(\|Au\|\right))\varphi(dA)\nu(du) \ < \ 0.
\end{eqnarray*}
The claim follows. \end{proof}


\begin{thebibliography}{99}
\bibitem[Arnold]{Arnold} Arnold, L., 1998. \textit{Random Dynamical Systems}. Springer.
\bibitem[Benaim]{Benaim} Benaim, M.\ and Hurth, T., 2022. \textit{Markov Chains on Metric Spaces: A Short Course}. Springer, 2022.
\bibitem[Bill]{Bill} Billingsley, P., 1999. \textit{Convergence of probability measures} (2nd ed.), Wiley-Interscience.
\bibitem[FK]{FK} Furstenberg, H. and Kifer, Y., 1983. Random matrix products and measures on projective spaces. Israel Journal of Mathematics, 46(1), pp.12-32.
\bibitem[MeynTweedie]{MeynTweedie} Meyn, S.\ P.\ and Tweedie, R.\ L., \textit{Markov Chains and Stochastic Stability} (2nd ed.). Cambridge University Press, 2012.
\bibitem[Osel]{Osel} Oseledec, V.\ I., 1968. A multiplicative ergodic theorem, Liapunov characteristic numbers for dynamical systems. Transactions of the Moscow Mathematical Socierty 19, 197-221.
\end{thebibliography}
\end{document}
